\begin{abstract}
This report documents a single pipeline for comparing data-availability practices
between quantum computing research and classical (``regular'') computing research:
whether a paper relies on data external to itself, on data it generates itself, or
both, and how completely and verifiably each was shared. Candidate papers for each
side are retrieved directly from the relevant arXiv subject categories (quantum
physics and computing categories for the quantum side; the corresponding
computer-science categories for the classical side) over a fixed publication window,
with no keyword filtering of any kind. An earlier, keyword-based retrieval approach
was tried first and found to systematically miss papers that describe relevant work
using different terminology than any fixed keyword list anticipates --- a structural
limitation no amount of list-tuning can fully close --- so retrieval was switched to
this comprehensive, category-based pull instead. The result is a complete candidate
corpus of \textbf{109{,}321 quantum papers and 345{,}107 classical-computing papers}.
Each candidate is screened for data/code-availability signal via a full-text scan and
classified by a two-axis, six-case framework: external-data use and self-generated-data
use are scored independently, each across six states ranging from directly accessible
to claimed-but-unverifiable to not present at all. Because the comprehensive pull is
not filtered by topic at retrieval time, a separate step also flags papers that are not
really quantum-computing work at all (e.g.\ quantum-hardware papers about sensing
rather than computation), so that the pull's breadth does not come at the cost of
off-topic material in the final results. At the time of writing, 3{,}472 quantum
papers have been classified under this framework, drawn from an earlier, smaller,
keyword-curated subset built before retrieval was switched to full category-based
coverage; the classical side has not yet been classified. Classifying the remainder of
the now much larger corpus is a substantially bigger undertaking than the work
completed so far, and is the natural next phase of this project.
\end{abstract}
